\documentclass[12pt]{article}
\usepackage[utf8]{inputenc}
\usepackage[T1]{fontenc}
\usepackage[margin=1in]{geometry}
\usepackage{setspace}
\usepackage{xcolor}
\usepackage[colorlinks=true, urlcolor=blue!60!black, linkcolor=blue!60!black]{hyperref}
\setlength{\parindent}{0em}
\setlength{\parskip}{0.45em}
\renewcommand{\baselinestretch}{1.0}
\begin{document}
\begin{center}
{\Large \textit{Poindexter v. Greenhow}, 114 U.S. 270 (1885)\par}
\vspace{0.3em}
{\normalsize Case Brief --- One-Page Edition\par}
\end{center}

\textbf{Memory Jogger.} The lead Virginia coupon case: an officer enforcing a void statute is not the State (\href{https://supreme.justia.com/cases/federal/us/114/270/}{114 U.S. 270}).

\textbf{Facts.} Virginia's 1871 bonds carried coupons receivable for taxes; an 1882 Act forbade collectors to accept them. Richmond treasurer Samuel Greenhow refused the matured coupons Poindexter tendered for \$12.45 in taxes and seized his office desk under a distraint statute.

\textbf{Procedural History.} Poindexter's detinue suit against Greenhow failed in the Hustings Court of Richmond.

\textbf{Votes.} Reversed, 5--4. Matthews wrote for the Court; Bradley, Waite, Miller, and Gray dissented.

\textbf{Holding.} The coupons are not bills of credit; the 1882 Act impaired the obligation of contracts and was void; and the collector could be sued personally despite the Eleventh Amendment.

\textbf{Reasoning.} The coupons were receivable only at maturity, not designed to circulate as money. A lawful tender equals payment, so Greenhow lacked authority, and a void statute is ``as if it were not and never had been.'' Suits to compel payment from the treasury remain barred. Bradley called the distinction a fiction; \textit{Ex parte Young} later adopted it.

\textbf{Examples.} \emph{Same-Side:} suing an officer who seized a truck after refusing coupons. \emph{Opposite-Side:} ordering the treasurer to pay coupons. \emph{Fence-Sitter:} an injunction to accept coupons.

\textbf{Critique.} Greenhow is a state actor for the merits but a private wrongdoer for immunity, a paradox that remains the defining fiction of federal courts doctrine.

\textbf{Key Quotation.} ``The defendant in error is not her officer, her agent, or her representative, in the matter complained of; for he has acted not only without her authority, but contrary to her express commands.'' --- Justice Matthews, Opinion of the Court, 114 U.S. at 292--93.

\end{document}
